\subsection{Rational input bit-length complexity}\label{sec:bits}

\begin{proposition}[Required order in the input size]\label{prop:bits}
Encode each of the eight rational atoms by its reduced numerator and positive denominator in ordinary binary, with delimiters of linear overhead. For a law $P$ on three bits, let $\ell(P)$ be its total bit length. For each $H\in\{\mathrm{NW},\mathrm{AI},\mathrm{exp}\}$,
\[
\max\bigl\{\tmin^H(P):\ \ell(P)\le B,\ P\notin\Ctri\bigr\}=2^{\Theta(B)} .
\]
\end{proposition}

\begin{proof}[Proof of Proposition~\ref{prop:bits}]
\emph{Lower bound.} With $k=2^j$, $\eps=k^{-3}$ and $\sigma=1/(2k^2)$, the eight atoms of $P_\eps$ are rationals with denominator $8k^9$ and numerators of $O(j)$ bits, so $\ell(P_\eps)=O(j)$, while Proposition~\ref{prop:family} gives $\tmin^H(P_\eps)>k/2=2^{j-1}$. Choosing $j$ proportional to $B$ gives $2^{\Omega(B)}$.

\emph{Upper bound.} By \cite{RGW2018}, $\Ctri$ is the image of a compact semialgebraic parameter space (six latent values per source and stochastic binary response tables) under a polynomial map, hence a closed semialgebraic subset of the simplex. The function $u(P)=d_2(P,\Ctri)^2$ is therefore semialgebraic, and its graph is described by a fixed quantifier-free formula $\Phi(P,u)$ in integer polynomials of fixed degrees, by real quantifier elimination \cite{BasuPollackRoy2006}. For a rational $P$ of bit length $B$, substituting $P$ and clearing denominators turns the atoms of $\Phi$ into univariate integer polynomials in $u$ of fixed degree and coefficient height $2^{O(B)}$. The set of $u$ satisfying $\Phi(P,u)$ is the single point $u(P)$; discard identically zero specialized polynomials, whose comparisons have constant truth values. If none of the remaining polynomials vanished at $u(P)$, their signs would persist on a neighbourhood and the formula would define an interval. Thus some nonzero specialized polynomial $f$ vanishes at $u(P)$. Remove any factor $u^r$ from $f$ to obtain $g(u)=a_0+\cdots+a_du^d$ with $a_0\ne0$, the same nonzero root, and coefficient height at most $H=2^{O(B)}$. For $0<u<1/(H+1)$,
\[
\Bigl|\sum_{i=1}^d a_i u^i\Bigr|\le H\sum_{i=1}^d u^i<\frac{Hu}{1-u}<1\le|a_0|,
\]
so $g(u)\ne0$. Hence $u(P)\ge2^{-O(B)}$ when $P$ is incompatible, and Corollary~\ref{prop:promised} gives $\tmin^{\mathrm{NW}}(P)\le(1-\|P\|_2^2)/u(P)+1=2^{O(B)}$; the two smaller feasible sets reject no later.
\end{proof}

With the binary triangle held fixed, a quantifier-free membership formula contains a bounded number of integer polynomials of bounded degree. Evaluating their signs at a rational input takes time polynomial in $B$. Quantifier elimination supplies such a formula in principle \cite{BasuPollackRoy2006}, although deriving and storing it may be impractical. For this fixed scenario, a polynomial-time decision procedure therefore exists alongside the exponential bound on the inflation order. For an incompatible rational input, one can also compute the first rejecting order by enumerating exact rational feasibility problems.
